% Vietnamese translation for OI Public Library
% Source-File: IOI中国国家候选队论文/2005/周源/周源.ppt
\documentclass[11pt]{article}
\usepackage{oipl-vn}

\newcommand{\Oh}{\mathcal{O}}
\newcommand{\lei}{\ensuremath{\mathrm{lei}}}
\newcommand{\euro}{\ensuremath{\mathrm{euro}}}

\title{Bản trình chiếu: Phương pháp nén trong thi tin học}
\author{Zhou Yuan}
\date{2005}

\begin{document}
\maketitle

\origtitle{Bỏ phần thừa, giữ tinh hoa: phương pháp nén trong thi tin học, bản trình chiếu}
\sourcepath{IOI China candidate team papers / 2005 / Zhou Yuan / Zhou Yuan.ppt}

\section{Mạch chính của slide}

Bản trình chiếu đi cùng bài viết \emph{Bỏ phần thừa, giữ tinh hoa}. Phần chữ
khôi phục được từ PowerPoint cũ nhắc tới \textit{BOI 2003 Euro}, \textit{ZOJ
p1794}, các dãy
\[
  a_{1,1},a_{1,2},\ldots,a_{1,n_1},\quad
  a_{2,1},\ldots,a_{2,n_2},\quad
  \ldots,\quad
  a_{K,1},\ldots,a_{K,n_K},
\]
các mốc độ phức tạp \(\Oh(KNK)\), \(\Oh(N\log^2K)\), và các công thức đồng
dư. Bản ghi chú này dùng bài viết đã dịch làm neo để phục hồi mạch slide:
định nghĩa phương pháp nén, hai yếu tố then chốt, rồi hai ứng dụng chính.

\section{Định nghĩa của phương pháp nén}

Trong bài viết, ``nén'' không chỉ là giảm dung lượng như tệp nén. Nó là một
cách giải bài toán: bỏ thông tin dư thừa, giữ phần tinh hoa đủ ảnh hưởng tới
đáp án, rồi thay một khối dữ liệu phức tạp bằng một đối tượng đơn giản hơn.

Giả sử có một tập \(S\). Nếu quan hệ giữa các phần tử bên trong \(S\) không
còn ảnh hưởng tới kết quả, và phần còn lại của bài toán chỉ cần biết cách
\(S\) liên hệ với bên ngoài, ta có thể xem \(S\) là một gói. Khi nén, ta bỏ
chi tiết nội bộ của gói, giữ các liên hệ cần thiết, rồi dùng một phần tử mới
thay cho cả gói.

Hai ví dụ đơn giản là nhiều nguồn trong đường đi ngắn nhất và đội bóng truyền
tin. Với nhiều nguồn, ta nén cả tập nguồn thành một siêu nguồn để chạy
Dijkstra một lần. Với đội bóng, mỗi thành phần liên thông mạnh có trạng thái
biết tin giống nhau; nén các thành phần ấy thành đỉnh của một DAG, đáp án là
số đỉnh có bậc vào bằng \(0\).

\section{Hai yếu tố then chốt}

Yếu tố thứ nhất là \emph{tính nén được}. Ta phải biết vì sao thông tin nội bộ
có thể bỏ. Nếu bỏ nhầm quan hệ còn ảnh hưởng tới đáp án, lời giải sai dù chạy
rất nhanh. Trong ví dụ nhiều nguồn, đi lại giữa hai nguồn không giúp đường đi
ngắn nhất tốt hơn, nên các cạnh nội bộ của tập nguồn là thừa. Trong ví dụ
thành phần liên thông mạnh, mọi đỉnh trong cùng thành phần có thể thông báo
cho nhau, nên chỉ cần biết có thành phần ấy hay không.

Yếu tố thứ hai là \emph{quy tắc thay thế}. Nén không phải là xóa sạch; phải
giữ lại đúng thông tin bên ngoài. Với đồ thị, quy tắc có thể là thay một tập
đỉnh bằng một đỉnh mới và nối các cạnh ra ngoài. Với số học, quy tắc có thể
là thay hai điều kiện bằng một điều kiện tương đương.

Slide có các công thức đồng dư:
\[
  a_1x \equiv c_1 \pmod {b_1},\qquad
  a_2x \equiv c_2 \pmod {b_2}.
\]
Mỗi phương trình có thể đưa về dạng
\[
  x=x_1+p_1t,\qquad x=x_2+p_2s.
\]
Nếu hai tập nghiệm giao nhau, giao ấy lại là một lớp đồng dư
\[
  x\equiv x_0 \pmod{\operatorname{lcm}(p_1,p_2)}.
\]
Đó là một phép nén: hai ràng buộc được thay bằng một ràng buộc giữ nguyên tập
nghiệm.

\section{\textit{Euro}: nén để khôi phục tính đơn điệu}

\textit{Euro} của BOI 2003 cho dãy thu nhập euro \(a_1,\ldots,a_N\), mỗi
ngày \(i\) có tỷ giá \(1\euro:i\lei\). Mỗi lần đổi phải trả phí cố định
\(T\lei\), và tới ngày \(N\) phải đổi hết. Cần tối đa hóa số \(\lei\) nhận
được.

Nếu chọn các ngày đổi \(c_1<c_2<\cdots<c_k=N\), lợi nhuận là tổng các khoản
thu trong từng đoạn, nhân với ngày đổi của đoạn, trừ phí. Quy hoạch động tự
nhiên là
\[
  f(n)=\max_{0\le i<n}\{f(i)+n(S_n-S_i)-T\},
\]
trong đó \(S_i\) là tổng tiền tố. Công thức này giống dạng tối ưu hóa bao
lồi, nhưng \(a_i\) có thể âm nên \(S_i\) không đơn điệu. Đây là mâu thuẫn
chính của ví dụ.

Nhận xét nén là: nếu từ ngày \(i\), các tổng tiền tố của đoạn vẫn dương cho
tới trước ngày \(j\), còn tổng cả đoạn tới \(j\) không dương, thì trong một
phương án tối ưu không cần đổi giữa chừng trong đoạn ấy. Nếu đang giữ euro
dương, đổi muộn hơn tốt hơn; nếu đang giữ euro âm, có thể dời điểm đổi về
trước đoạn mà không làm xấu kết quả. Vì vậy cả đoạn có thể thay bằng một cặp
\((b,t)\): \(b\) là tổng euro của đoạn, \(t\) là ngày kết thúc đoạn.

Sau khi nén liên tiếp, các khối chính có tổng không dương, nên dãy tổng tiền
tố mới trở nên đơn điệu theo hướng cần thiết. Điều kiện của tối ưu hóa bao
lồi được phục hồi. Thuật toán vì thế gồm một lần quét để nén và một lần quy
hoạch động tối ưu hóa bằng bao lồi, tổng thời gian \(\Oh(N)\). Đây là ví dụ
đẹp vì nén không chỉ giảm số phần tử; nó còn biến dữ liệu về dạng có tính
đơn điệu.

\section{\textit{Merging Sequence}: nén hai giai đoạn}

Ví dụ lớn thứ hai là \textit{ZOJ p1794 Merging Sequence Problem}. Có \(K\)
dãy, độ dài \(n_1,\ldots,n_K\), tổng độ dài \(N\). Mỗi bước chọn phần tử đầu
của một dãy còn lại và đưa vào dãy kết quả \(S\). Mục tiêu là làm cho giá trị
lớn nhất của các tổng tiền tố trong \(S\) nhỏ nhất có thể.

Nếu xử lý trực tiếp, trạng thái phải nhớ đã lấy bao nhiêu phần tử từ mỗi dãy,
không khả thi. Ta tìm các đoạn trong cùng một dãy mà khi đã bắt đầu lấy thì
có thể lấy liên tiếp trong một nghiệm tối ưu. Đoạn như vậy được nén thành một
cặp \((a,b)\): \(a\ge0\) là đỉnh tương đối lớn nhất bên trong đoạn, còn
\(a+b\) là tổng của cả đoạn, với \(b\le0\). Cặp này mô tả đủ ảnh hưởng của
đoạn lên tổng tiền tố toàn cục.

Giai đoạn nén thứ nhất xét một đoạn có các tổng tiền tố trước cuối đều không
âm, còn tổng cả đoạn âm. Khi đã lấy phần tử đầu của đoạn, hoặc chưa lấy gì
thêm, hoặc lấy liên tiếp tới hết đoạn; xen phần tử khác vào giữa không giúp
giảm đỉnh lớn nhất. Vì vậy đoạn được nén thành một cặp âm.

Giai đoạn nén thứ hai xử lý các cặp âm liên tiếp. Nếu đỉnh của hai cặp nằm ở
cặp thứ nhất, tức
\[
  a_1\ge a_1+b_1+a_2,
\]
thì đưa cặp thứ hai đứng ngay sau cặp thứ nhất không làm tăng đỉnh lớn nhất.
Do đó hai cặp có thể nén tiếp. Lặp lại điều kiện này làm cho trong mỗi dãy,
đỉnh tương đối của các cặp còn lại tăng dần.

\section{Tham lam sau khi nén}

Sau hai giai đoạn, bài toán còn \(K\) dãy các cặp có tổng âm và đỉnh tương
đối tăng dần trong từng dãy. Vì mỗi cặp làm tổng hiện tại giảm sau khi đi qua,
cặp có đỉnh nhỏ nên được xử lý sớm; cặp có đỉnh lớn nên để muộn hơn khi tổng
nền đã thấp hơn. Do từng dãy đã có thứ tự tăng, lời giải là trộn \(K\) dãy
đã sắp xếp theo đỉnh tương đối.

Dùng heap lưu cặp đầu của mỗi dãy, mỗi lần lấy cặp có đỉnh nhỏ nhất rồi đẩy
cặp kế tiếp của dãy đó. Thời gian là
\[
  \Oh(N\log K).
\]
Phần mã trong tài liệu gốc chỉ triển khai hai giai đoạn nén và phép trộn bằng
heap; bản ghi chú không chép lại chương trình.

\section{Kết luận}

Phương pháp nén gồm hai câu hỏi: có thể bỏ thông tin nội bộ nào, và thay phần
bị bỏ bằng đại diện nào. Nếu trả lời đúng, bài toán không chỉ nhỏ hơn mà còn
có cấu trúc rõ hơn: nhiều nguồn thành một nguồn, thành phần liên thông mạnh
thành DAG, dãy \textit{Euro} thành dãy đơn điệu, và \(K\) dãy lộn xộn thành
các cặp có thể trộn tham lam.

Vì vậy khẩu hiệu của slide là ``bỏ phần thừa, giữ tinh hoa''. Nén không phải
mẹo cài đặt cuối cùng; nó là một bước mô hình hóa thông tin trước khi chọn
thuật toán.

\end{document}
